\section{طول یک مدول}

% Source: sources/lectures/13 Length of a Module.pdf | PDF page: 1 | item: series and refinement
\begin{frame}[t]{سری و تظریف سری}
  \label{lec13:series-refinement}
  \begin{definition}[سری]
    زنجیری از زیرمدول‌های $M$ به شکل
    \[
      0=M_0\leqslant M_1\leqslant\cdots\leqslant M_n=M
    \]
    را یک \textbf{سری به طول $n$} می‌نامیم.
  \end{definition}
  \begin{definition}[تظریف سری]
    سری $\{M_i\}_{i=1}^{n}$ را تظریفی از سری
    $\{M'_j\}_{j=1}^{m}$ گوییم هرگاه
    \[
      \{M'_j\}_{j=1}^{m}\subseteq\{M_i\}_{i=1}^{n}.
    \]
  \end{definition}
\end{frame}

% Source: sources/lectures/13 Length of a Module.pdf | PDF page: 1 | item: factors and composition series
\begin{frame}[t]{عامل‌های سری و سری ترکیبی}
  \label{lec13:composition-series}
  \begin{definition}[عامل سری]
    در سری
    \[
      0=M_0\leqslant M_1\leqslant\cdots\leqslant M_n=M,
    \]
    به هر یک از مدول‌های
    \[
      M_{i+1}/M_i\qquad(0\leqslant i<n)
    \]
    یک \textbf{عامل سری} گفته می‌شود.
  \end{definition}
  \begin{definition}[سری ترکیبی]
    اگر هر عامل سری یک $R$-مدول ساده باشد، آن سری را یک
    \textbf{سری ترکیبی} برای $M$ می‌نامیم.
  \end{definition}
\end{frame}

% Source: sources/lectures/13 Length of a Module.pdf | PDF page: 1 | item: finite direct sum of simple modules
\begin{frame}[t]{سری ترکیبیِ جمع مستقیم مدول‌های ساده}
  \label{lec13:simple-sum-example}
  \begin{example}
    اگر $\{M_i\}_{i=1}^{n}$ خانواده‌ای از $R$-مدول‌های ساده باشد، آنگاه
    \[
      0\leqslant M_1\leqslant M_1\oplus M_2
      \leqslant\cdots\leqslant\bigoplus_{i=1}^{n}M_i
    \]
    یک سری ترکیبی برای
    \[
      \bigoplus_{i=1}^{n}M_i
    \]
    است.
  \end{example}
\end{frame}

% Source: sources/lectures/13 Length of a Module.pdf | PDF page: 1 | item: Z9 example and Z nonexample
\begin{frame}[t]{دو مثال از مدول‌های روی اعداد صحیح}
  \label{lec13:integer-examples}
  \begin{example}
    زنجیر
    \[
      0\lneq\{\overline{0},\overline{3},\overline{6}\}
      \lneq\mathbb Z_9
    \]
    یک سری ترکیبی برای $\mathbb Z_9$ به‌عنوان $\mathbb Z$-مدول است.
  \end{example}
  \begin{example}
    $\mathbb Z$ به‌عنوان $\mathbb Z$-مدول سری ترکیبی ندارد.
    \hfill (چرا؟)
  \end{example}
\end{frame}

% Source: sources/lectures/13 Length of a Module.pdf | PDF page: 2 | item: length definition
\begin{frame}[t]{طول یک مدول}
  \label{lec13:length-definition}
  \begin{definition}[طول مدول]
    فرض کنید $M$ یک $R$-مدول باشد که حداقل یک سری ترکیبی دارد.
    طول کوتاه‌ترین سری ترکیبیِ $M$ را با $\ell(M)$ نمایش می‌دهیم و آن را
    \textbf{طول مدول $M$} می‌گوییم.
    \par\medskip
    اگر سری ترکیبی برای $M$ وجود نداشته باشد، تعریف می‌کنیم
    \[
      \ell(M)=\infty.
    \]
  \end{definition}
  در حالت اول، $M$ مدولی با طول متناهی نامیده می‌شود.
\end{frame}

% Source: sources/lectures/13 Length of a Module.pdf | PDF page: 2 | item: proper submodule length theorem and intersection chain
\begin{frame}[t]{طول زیرمدول سره}
  \label{lec13:proper-submodule-length}
  \begin{theorem}
    اگر $M$ یک $R$-مدول با طول متناهی و $N$ زیرمدولی سره از $M$ باشد، آنگاه
    \[
      \ell(N)<\ell(M).
    \]
  \end{theorem}
  \textbf{آغاز اثبات.}
  فرض کنید $\ell(M)=n$ و
  \[
    0=M_0\leqslant M_1\leqslant\cdots\leqslant M_n=M
  \]
  یک سری ترکیبی به طول $n$ باشد. با اشتراک گرفتن با $N$ داریم
  \[
    0=M_0\cap N\leqslant M_1\cap N
    \leqslant\cdots\leqslant M_n\cap N=N.
  \]
  ممکن است بعضی از جمله‌های این زنجیر برابر باشند.
\end{frame}

% Source: sources/lectures/13 Length of a Module.pdf | PDF pages: 2-3 | item: canonical homomorphism and possible images
\begin{frame}[t]{عامل‌های زنجیر اشتراک‌ها}
  \label{lec13:intersection-factors}
  برای $1\leqslant i\leqslant n$ تعریف می‌کنیم
  \[
    \varphi_i:M_i\cap N\longrightarrow M_i/M_{i-1},
    \qquad x\longmapsto x+M_{i-1}.
  \]
  بنابر قضیهٔ اول یکریختی،
  \[
    \ker\varphi_i=M_{i-1}\cap N,
    \qquad
    \frac{M_i\cap N}{M_{i-1}\cap N}
      \cong\operatorname{Im}\varphi_i\leqslant M_i/M_{i-1}.
  \]
  چون $M_i/M_{i-1}$ ساده است، یا $\operatorname{Im}\varphi_i=0$، یا
  \[
    \operatorname{Im}\varphi_i=M_i/M_{i-1}.
  \]
  در حالت دوم، عامل متناظر ساده است. بنابراین $\ell(N)\leqslant\ell(M)$.
\end{frame}

% Source: sources/lectures/13 Length of a Module.pdf | PDF page: 3 | item: assume all images nonzero and induction setup
\begin{frame}[t]{وجود یک تصویر صفر}
  \label{lec13:zero-image-induction}
  \textbf{ادعا.} دست‌کم یکی از تصویرهای $\operatorname{Im}\varphi_i$ صفر است.
  \par\medskip
  در غیر این صورت، همهٔ نگاشت‌های $\varphi_i$ روی عامل‌های ساده پوشا هستند.
  با استقرا نشان می‌دهیم
  \[
    M_i\cap N=M_i\qquad(0\leqslant i\leqslant n).
  \]
  پایهٔ استقرا، $M_0\cap N=M_0=0$، واضح است. \hfill(چرا؟)
  \par\medskip
  فرض کنید $M_{i-1}\cap N=M_{i-1}$. آنگاه، با همان نگاشتِ طبیعی،
  \[
    \frac{M_i\cap N}{M_{i-1}\cap N}
      \cong\frac{M_i}{M_{i-1}},
    \qquad M_i\cap N\subseteq M_i.
  \]
\end{frame}

% Source: sources/lectures/13 Length of a Module.pdf | PDF pages: 3-4 | item: induction element calculation and contradiction
\begin{frame}[t]{پایان اثبات نامساویِ طول}
  \label{lec13:proper-length-proof-end}
  اگر $x\in M_i$، از پوشایی $\varphi_i$ عنصری مانند $t\in M_i\cap N$ هست که
  \[
    x+M_{i-1}=t+M_{i-1}.
  \]
  پس با فرض استقرا،
  \[
    x-t\in M_{i-1}=M_{i-1}\cap N.
  \]
  از $t\in N$ نتیجه می‌شود $x\in N$ و در نتیجه $x\in M_i\cap N$.
  بنابراین $M_i\subseteq M_i\cap N$ و لذا $M_i=M_i\cap N$.
  \par\medskip
  در انتهای استقرا،
  \[
    M=M_n=M_n\cap N=M\cap N=N,
  \]
  که با سره بودن $N$ تناقض دارد. پس یک تصویر صفر است و
  \[
    \ell(N)<\ell(M).\qed
  \]
\end{frame}

% Source: sources/lectures/13 Length of a Module.pdf | PDF page: 4 | item: Jordan-Holder statement as given (length and refinement only)
\begin{frame}[t]{قضیهٔ جردن هولدر}
  \label{lec13:jordan-holder}
  \begin{theorem}[جردن هولدر]
    اگر $M$ یک $R$-مدول با طول $n$ باشد، هر سری ترکیبیِ $M$ طول $n$ دارد.
    \par\medskip
    به‌علاوه، هر سری از زیرمدول‌های $M$ تظریفی دارد که یک سری ترکیبیِ $M$ است.
  \end{theorem}
\end{frame}

% Source: sources/lectures/13 Length of a Module.pdf | PDF page: 4 | item: bound on strict series and equality of composition lengths
\begin{frame}[t]{طول سری‌های ترکیبی}
  \label{lec13:composition-length-proof}
  \textbf{اثبات.} داریم $\ell(M)=n$. طول هر سریِ سره از زیرمدول‌های $M$
  حداکثر $n$ است، زیرا در زنجیر
  \[
    0\lneq M_1\lneq M_2\lneq\cdots\lneq M_k=M
  \]
  با نتیجهٔ قبل،
  \[
    0=\ell(0)<\ell(M_1)<\ell(M_2)<\cdots<\ell(M_k)=n.
  \]
  در هر گام، طول دست‌کم یک واحد افزایش می‌یابد. پس $k\leqslant n$.
  \par\medskip
  اگر سری ترکیبی باشد، از تعریفِ طولِ کوتاه‌ترین سری ترکیبی $n\leqslant k$.
  بنابراین
  \[
    k=n=\ell(M).
  \]
\end{frame}

% Source: sources/lectures/13 Length of a Module.pdf | PDF page: 5 | item: refine a nonsimple factor
\begin{frame}[t]{تظریف یک سری به سری ترکیبی}
  \label{lec13:refinement-proof}
  اگر عامل $M_i/M_{i-1}$ ساده نباشد، زیرمدولی دارد که
  \[
    0\neq N/M_{i-1}\lneq M_i/M_{i-1}.
  \]
  در نتیجه،
  \[
    M_{i-1}\lneq N\lneq M_i.
  \]
  با افزودن $N$ به سری، طولِ سری یک واحد افزایش می‌یابد.
  \par\medskip
  از طرفی طول سری‌های سره کرانِ $n$ دارد. بنابراین این فرایند سرانجام
  به یک سری ترکیبی می‌رسد.\hfill$\qed$
\end{frame}

% Source: sources/lectures/13 Length of a Module.pdf | PDF page: 5 | item: isomorphic powers of simple modules and canonical series
\begin{frame}[t]{توان‌های یکریختِ مدول‌های ساده}
  \label{lec13:simple-powers}
  \begin{corollary}
    اگر $M$ و $N$ دو $R$-مدول ساده باشند و
    \[
      M^m\cong N^n,
    \]
    آنگاه $m=n$ و $M\cong N$.
  \end{corollary}
  \textbf{آغاز اثبات.} سری ترکیبیِ زیر برای $M^m$ را در نظر بگیرید:
  \[
    0\lneq M\lneq M\oplus M\lneq\cdots\lneq M^m.
  \]
  اگر $\varphi:M^m\longrightarrow N^n$ یک $R$-یکریختی باشد، تصویر این زنجیر
  سری‌ای برای $N^n$ می‌دهد.
\end{frame}

% Source: sources/lectures/13 Length of a Module.pdf | PDF page: 5 | item: transported factors remain simple and equal powers
\begin{frame}[t]{سادگیِ عامل‌های سریِ منتقل‌شده}
  \label{lec13:transported-series}
  زنجیرِ تصویرها عبارت است از
  \[
    0\lneq\varphi(M)\lneq\varphi(M\oplus M)
      \lneq\cdots\lneq\varphi(M^m)=N^n.
  \]
  هر عامل این سری ساده است. مثلاً اگر
  \[
    \varphi(M)\lneq T'\lneq\varphi(M\oplus M),
  \]
  آنگاه
  \[
    M\lneq\varphi^{-1}(T')\lneq M\oplus M,
  \]
  که با سادگیِ عامل متناظر تناقض دارد.
  \par\medskip
  پس این زنجیر، سری ترکیبیِ $N^n$ به طول $m$ است. سریِ طبیعیِ
  $0\lneq N\lneq N\oplus N\lneq\cdots\lneq N^n$ نیز طول $n$ دارد.
  بنابراین $m=n$.
\end{frame}

% Source: sources/lectures/13 Length of a Module.pdf | PDF page: 5 | item: Hom decomposition and Schur-lemma contradiction left as why in source
\begin{frame}[t]{یکریختیِ دو مدول ساده}
  \label{lec13:simple-powers-hom}
  از یکریختیِ داده‌شده و تجزیهٔ همریختی‌های جمع مستقیم داریم
  \[
    \operatorname{Hom}_R(N^n,N)
      \cong\operatorname{Hom}_R(M^m,N)
      \cong\bigoplus_{i=1}^{m}\operatorname{Hom}_R(M,N).
  \]
  این یکریختی‌ها، یکریختیِ $\mathbb Z$-مدولی‌اند.
  \par\medskip
  اگر $M\not\cong N$، بنابر لم شور به تناقض می‌رسیم. \hfill(چرا؟)
  \par\medskip
  در نتیجه $M\cong N$.\hfill$\qed$
\end{frame}

% Source: sources/lectures/13 Length of a Module.pdf | PDF page: 6 | item: finite length iff Noetherian and Artinian; forward implication
\begin{frame}[t]{طول متناهی و شرط‌های زنجیری}
  \label{lec13:finite-length-chain-conditions}
  \begin{theorem}
    برای یک $R$-مدول $M$،
    \[
      \ell(M)<\infty
      \quad\Longleftrightarrow\quad
      M\text{ نوتری و آرتینی است.}
    \]
  \end{theorem}
  \textbf{اثباتِ جهت اول.}
  اگر $\ell(M)<\infty$، طول سری‌های سره کران متناهی دارد.
  بنابراین زنجیرهای صعودی و نزولیِ زیرمدول‌ها متوقف می‌شوند و $M$ نوتری و آرتینی است.
\end{frame}

% Source: sources/lectures/13 Length of a Module.pdf | PDF page: 6 | item: construct minimal submodules successively in quotients
\begin{frame}[t]{ساخت سری ترکیبی از شرط‌های زنجیری}
  \label{lec13:chain-conditions-converse}
  \textbf{عکس.} فرض کنید $M$ نوتری و آرتینی باشد.
  از آرتینی بودن، یک زیرمدول مینیمال $M_1$ داریم که ساده است.
  اگر $M=M_1$، ساخت پایان یافته است.
  \par\medskip
  در غیر این صورت، $M/M_1$ آرتینی است و زیرمدولی مینیمال مانند
  $M_2/M_1$ دارد که ساده است. اگر $M=M_2$، ساخت پایان یافته است؛
  وگرنه همین فرایند را ادامه می‌دهیم.
  \[
    0\lneq M_1\lneq M_2\lneq M_3\lneq\cdots
  \]
  چون $M$ نوتری است، زنجیر متوقف می‌شود. پس برای یک $i$، $M_i=M$ و
  سری ترکیبی‌ای به طول $i$ داریم. در نتیجه $\ell(M)<\infty$.\hfill$\qed$
\end{frame}

% Source: sources/lectures/13 Length of a Module.pdf | PDF page: 6 | item: finite length of submodules and quotients
\begin{frame}[t]{طول متناهیِ زیرمدول و مدول خارج‌قسمتی}
  \label{lec13:finite-length-subquotients}
  \begin{corollary}
    فرض کنید $M$ یک $R$-مدول با طول متناهی و $N\leqslant M$ باشد.
    در این صورت $N$ و $M/N$ نیز طول متناهی دارند.
  \end{corollary}
  \textbf{اثبات.} از نتیجهٔ قبل، $M$ نوتری و آرتینی است.
  این دو خاصیت به $N$ و $M/N$ منتقل می‌شوند. بنابراین
  \[
    \ell(N)<\infty,
    \qquad
    \ell(M/N)<\infty.
  \]
  \hfill$\qed$
\end{frame}

% Source: sources/lectures/13 Length of a Module.pdf | PDF page: 6 | item: image/kernel eventual sum and intersection theorem
\begin{frame}[t]{تصویر و هستهٔ توان‌های خودهمریختی}
  \label{lec13:eventual-image-kernel}
  \begin{theorem}
    فرض کنید $M$ یک $R$-مدول و $f:M\longrightarrow M$ یک $R$-همریختی باشد.
    \par\medskip
    \textbf{الف)} اگر $M$ آرتینی باشد، عددی مانند $a$ هست که
    \[
      \operatorname{Im}f^a+\ker f^a=M.
    \]
    \textbf{ب)} اگر $M$ نوتری باشد، عددی مانند $b$ هست که
    \[
      \operatorname{Im}f^b\cap\ker f^b=0.
    \]
  \end{theorem}
\end{frame}

% Source: sources/lectures/13 Length of a Module.pdf | PDF page: 7 | item: stabilization of descending image chain (index renamed a)
\begin{frame}[t]{توقف زنجیر تصویرها}
  \label{lec13:image-stabilization}
  \textbf{اثباتِ الف.} زنجیر نزولی
  \[
    \operatorname{Im}f\supseteq\operatorname{Im}f^2
      \supseteq\operatorname{Im}f^3\supseteq\cdots
  \]
  را در نظر بگیرید.
  چون $M$ آرتینی است، عددی مانند $a$ هست که
  \[
    \operatorname{Im}f^a=\operatorname{Im}f^{a+i}
    \qquad(i\geqslant0).
  \]
  به‌ویژه،
  \[
    \operatorname{Im}f^a=\operatorname{Im}f^{2a}.
  \]
\end{frame}

% Source: sources/lectures/13 Length of a Module.pdf | PDF page: 7 | item: element decomposition for the sum (index renamed a)
\begin{frame}[t]{جمعِ تصویر و هسته}
  \label{lec13:image-kernel-sum}
  اگر $x\in M$، آنگاه
  \[
    f^a(x)\in\operatorname{Im}f^a=\operatorname{Im}f^{2a}.
  \]
  پس $y\in M$ هست که
  \[
    f^a(x)=f^{2a}(y)=f^a\bigl(f^a(y)\bigr).
  \]
  بنابراین $x-f^a(y)\in\ker f^a$ و
  \[
    x=\underbrace{f^a(y)}_{\in\operatorname{Im}f^a}
      +\underbrace{\bigl(x-f^a(y)\bigr)}_{\in\ker f^a}.
  \]
  نتیجه می‌شود
  \[
    M=\operatorname{Im}f^a+\ker f^a.\qed
  \]
\end{frame}

% Source: sources/lectures/13 Length of a Module.pdf | PDF page: 7 | item: stabilization of ascending kernel chain (index renamed b)
\begin{frame}[t]{توقف زنجیر هسته‌ها}
  \label{lec13:kernel-stabilization}
  \textbf{اثباتِ ب.} به‌طور مشابه، زنجیر صعودی
  \[
    \ker f\subseteq\ker f^2\subseteq\ker f^3\subseteq\cdots
  \]
  را داریم. \hfill(چرا؟)
  \par\medskip
  چون $M$ نوتری است، عددی مانند $b$ هست که
  \[
    \ker f^b=\ker f^{b+i}\qquad(i\geqslant0).
  \]
  به‌ویژه،
  \[
    \ker f^b=\ker f^{2b}.
  \]
\end{frame}

% Source: sources/lectures/13 Length of a Module.pdf | PDF page: 7 | item: zero intersection element argument (index renamed b)
\begin{frame}[t]{اشتراکِ تصویر و هسته}
  \label{lec13:image-kernel-intersection}
  اگر $x\in\ker f^b\cap\operatorname{Im}f^b$، آنگاه عنصری مانند $y\in M$ هست که
  \[
    x=f^b(y),\qquad f^b(x)=0.
  \]
  در نتیجه،
  \[
    f^{2b}(y)=f^b\bigl(f^b(y)\bigr)=f^b(x)=0.
  \]
  پس $y\in\ker f^{2b}=\ker f^b$ و بنابراین
  \[
    x=f^b(y)=0.
  \]
  نتیجه می‌شود
  \[
    \ker f^b\cap\operatorname{Im}f^b=0.\qed
  \]
\end{frame}

% Source: sources/lectures/13 Length of a Module.pdf | PDF page: 8 | item: literal Fitting statement, provisional exact exponent
% AMBIGUITY[L13-Q01]: primary, PDF page 8; length n in statement versus eventual m in proof. Exact-length exponent is provisional; its proof is explicitly deferred for optional later review.
\begin{frame}[t]{لم فیتینگ؛ صورت موقتِ دست‌نویس}
  \label{lec13:fitting-statement}
  \begin{question}[\lr{L13-Q01}: ابهام در توان]
    در صورت دست‌نویس، اگر $\ell(M)=n$ و $f:M\to M$ یک $R$-همریختی باشد،
    نتیجه به شکل زیر نوشته شده است:
    \[
      M=\operatorname{Im}f^n\oplus\ker f^n.
    \]
    اثبات همان صفحه، یک توانِ به‌اندازهٔ کافی بزرگ را انتخاب می‌کند.
    کرانِ توان‌های توقف با $\ell(M)$ در آن تصریح نشده است.
    \par\medskip
    صورت با توان $n=\ell(M)$ موقت است. توضیحِ کرانِ توان‌ها به‌عنوان
    موردی باز برای بازبینی بعدی ثبت شده است.
  \end{question}
\end{frame}

% Source: sources/lectures/13 Length of a Module.pdf | PDF page: 8 | item: simultaneous stabilization proof (source indices n,k renamed a,b)
% AMBIGUITY[L13-Q01]: this frame records the eventual-power proof only; it does not identify m with length(M).
\begin{frame}[t]{توقف هم‌زمان و تجزیهٔ مستقیم}
  \label{lec13:fitting-stabilization}
  چون $\ell(M)<\infty$، مدول $M$ نوتری و آرتینی است.
  توان‌های توقف تصویر و هسته را به‌ترتیب $a$ و $b$ بنامید.
  از دو اثبات قبل،
  \[
    M=\operatorname{Im}f^a+\ker f^a,
    \qquad
    \operatorname{Im}f^b\cap\ker f^b=0.
  \]
  اگر $m\geqslant\max\{a,b\}$، آنگاه
  \[
    \operatorname{Im}f^m=\operatorname{Im}f^a,
    \qquad
    \ker f^m=\ker f^b.
  \]
  پس با توقف هم‌زمان دو زنجیر،
  \[
    M=\operatorname{Im}f^m\oplus\ker f^m.
  \]
  \textbf{پرسش بازِ \lr{L13-Q01}:} ارتباط این توان با $n=\ell(M)$.
\end{frame}
